\begin{definition}
\label{definition-module-finite-type}
Let $R$ be a ring. Let $M$ be an $R$-module.
\begin{enumerate}
\item We say $M$ is a {\it finite $R$-module}, or a {\it finitely generated
$R$-module} if there exist $n \in \mathbf{N}$ and $x_1, \ldots, x_n \in M$
such that every element of $M$ is an $R$-linear combination of the $x_i$.
Equivalently, this means there exists a surjection
$R^{\oplus n} \to M$ for some $n \in \mathbf{N}$.
\item We say $M$ is a {\it finitely presented $R$-module} or an
{\it $R$-module of finite presentation} if there exist integers
$n, m \in \mathbf{N}$ and an exact sequence
$$
R^{\oplus m} \longrightarrow R^{\oplus n} \longrightarrow M \longrightarrow 0
$$
\end{enumerate}
\end{definition}

\begin{definition}
\label{definition-finite-type}
Let $R \to S$ be a ring map.
\begin{enumerate}
\item We say $R \to S$ is of {\it finite type}, or that {\it $S$ is a finite
type $R$-algebra} if there exist an $n \in \mathbf{N}$ and a surjection
of $R$-algebras $R[x_1, \ldots, x_n] \to S$.
\item We say $R \to S$ is of {\it finite presentation} if there
exist integers $n, m \in \mathbf{N}$ and polynomials
$f_1, \ldots, f_m \in R[x_1, \ldots, x_n]$
and an isomorphism of $R$-algebras
$R[x_1, \ldots, x_n]/(f_1, \ldots, f_m) \cong S$.
\end{enumerate}
\end{definition}

\begin{definition}
\label{definition-finite-ring-map}
Let $\varphi : R \to S$ be a ring map. We say $\varphi : R \to S$ is
{\it finite} if $S$ is finite as an $R$-module.
\end{definition}

\begin{lemma}
\label{lemma-finite-finitely-presented-extension}
Let $R \to S$ be a finite and finitely presented ring map.
Let $M$ be an $S$-module.
Then $M$ is finitely presented as an $R$-module if and only if
$M$ is finitely presented as an $S$-module.
\end{lemma}

\begin{proof}
One of the implications follows from
Lemma \ref{lemma-finitely-presented-over-subring}.
To see the other assume that $M$ is finitely presented as an $S$-module.
Pick a presentation
$$
S^{\oplus m} \longrightarrow
S^{\oplus n} \longrightarrow
M \longrightarrow 0
$$
As $S$ is finite as an $R$-module, the kernel of
$S^{\oplus n} \to M$ is a finite $R$-module. Thus from
Lemma \ref{lemma-extension}
we see that it suffices to prove that $S$ is finitely presented as an
$R$-module.

\medskip\noindent
Pick $y_1, \ldots, y_n \in S$ such that $y_1, \ldots, y_n$ generate $S$
as an $R$-module. By Lemma \ref{lemma-characterize-integral-element}
each $y_i$ is integral over $R$. Choose monic polynomials
$P_i(x) \in R[x]$ with $P_i(y_i) = 0$. Consider the ring
$$
S' = R[x_1, \ldots, x_n]/(P_1(x_1), \ldots, P_n(x_n))
$$
Then we see that $S$ is of finite presentation as an $S'$-algebra
by Lemma \ref{lemma-compose-finite-type}. Since $S' \to S$ is surjective,
the kernel $J = \Ker(S' \to S)$ is finitely generated as an ideal by
Lemma \ref{lemma-finite-presentation-independent}. Hence $J$ is a finite
$S'$-module (immediate from the definitions).
Thus $S = \Coker(J \to S')$ is  of finite presentation as an $S'$-module
by Lemma \ref{lemma-extension}.
Hence, arguing as in the first paragraph, it suffices to show that $S'$ is
of finite presentation as an $R$-module. Actually, $S'$ is free as an
$R$-module with basis the monomials $x_1^{e_1} \ldots x_n^{e_n}$
for $0 \leq e_i < \deg(P_i)$. Namely, write $R \to S'$ as the composition
$$
R \to R[x_1]/(P_1(x_1)) \to R[x_1, x_2]/(P_1(x_1), P_2(x_2)) \to
\ldots \to S'
$$
This shows that the $i$th ring in this sequence is free as a module over the
$(i - 1)$st one with basis $1, x_i, \ldots, x_i^{\deg(P_i) - 1}$. The result
follows easily from this by induction. Some details omitted.
\end{proof}

\begin{lemma}
\label{lemma-codimension}
Let $k$ be a field.
Let $S' \to S$ be a surjection of finite type $k$-algebras.
Let $\mathfrak p \subset S$ be a prime ideal,
and let $\mathfrak p'$ be the corresponding prime ideal of $S'$.
Let $X = \Spec(S)$, resp.\ $X' = \Spec(S')$,
and let $x \in X$, resp. $x'\in X'$ be the point corresponding
to $\mathfrak p$, resp.\ $\mathfrak p'$.
Then
$$
\dim_{x'} X' - \dim_x X =
\text{height}(\mathfrak p') - \text{height}(\mathfrak p).
$$
\end{lemma}

\begin{proof}
Immediate from Lemma \ref{lemma-dimension-at-a-point-finite-type-field}.
\end{proof}

\begin{lemma}
\label{lemma-dimension-at-a-point-preserved-field-extension}
Let $k$ be a field.
Let $S$ be a finite type $k$-algebra.
Set $X = \Spec(S)$.
Let $K/k$ be a field extension.
Set $S_K = K \otimes_k S$, and $X_K = \Spec(S_K)$.
Let $\mathfrak q \subset S$ be a prime corresponding to $x \in X$
and let $\mathfrak q_K \subset S_K$ be a prime corresponding
to $x_K \in X_K$ lying over $\mathfrak q$.
Then $\dim_x X = \dim_{x_K} X_K$.
\end{lemma}

\begin{proof}
Choose a presentation $S = k[x_1, \ldots, x_n]/I$.
This gives a presentation
$K \otimes_k S = K[x_1, \ldots, x_n]/(K \otimes_k I)$.
Let $\mathfrak q_K' \subset K[x_1, \ldots, x_n]$,
resp.\ $\mathfrak q' \subset k[x_1, \ldots, x_n]$ be
the corresponding primes. Consider the following
commutative diagram of Noetherian local rings
$$
\xymatrix{
K[x_1, \ldots, x_n]_{\mathfrak q_K'} \ar[r] &
(K \otimes_k S)_{\mathfrak q_K} \\
k[x_1, \ldots, x_n]_{\mathfrak q'} \ar[r] \ar[u] &
S_{\mathfrak q} \ar[u]
}
$$
Both vertical arrows are flat because they are localizations of
the flat ring maps $S \to S_K$ and
$k[x_1, \ldots, x_n] \to K[x_1, \ldots, x_n]$.
Moreover, the vertical arrows have the same fibre rings.
Hence, we see from
Lemma \ref{lemma-dimension-base-fibre-equals-total} that
$\text{height}(\mathfrak q') - \text{height}(\mathfrak q)
= \text{height}(\mathfrak q_K') - \text{height}(\mathfrak q_K)$.
Denote by $x' \in X' = \Spec(k[x_1, \ldots, x_n])$
and $x'_K \in X'_K = \Spec(K[x_1, \ldots, x_n])$
the points corresponding to $\mathfrak q'$ and
$\mathfrak q_K'$. By Lemma \ref{lemma-codimension} and what we showed
above we have
\begin{eqnarray*}
n - \dim_x X & = & \dim_{x'} X' - \dim_x X \\
& = & \text{height}(\mathfrak q') - \text{height}(\mathfrak q) \\
& = & \text{height}(\mathfrak q_K') - \text{height}(\mathfrak q_K) \\
& = & \dim_{x'_K} X'_K - \dim_{x_K} X_K \\
& = & n - \dim_{x_K} X_K
\end{eqnarray*}
and the lemma follows.
\end{proof}

\begin{lemma}
\label{lemma-integral-closure-localize}
Integral closure commutes with localization: If $A \to B$ is a ring
map, and $S \subset A$ is a multiplicative subset, then the integral
closure of $S^{-1}A$ in $S^{-1}B$ is $S^{-1}B'$, where $B' \subset B$
is the integral closure of $A$ in $B$.
\end{lemma}

\begin{proof}
Since localization is exact we see that $S^{-1}B' \subset S^{-1}B$.
Suppose $x \in B'$ and $f \in S$. Then
$x^d + \sum_{i = 1, \ldots, d} a_i x^{d - i} = 0$
in $B$ for some $a_i \in A$. Hence also
$$
(x/f)^d + \sum\nolimits_{i = 1, \ldots, d} a_i/f^i (x/f)^{d - i} = 0
$$
in $S^{-1}B$. In this way we see that $S^{-1}B'$ is contained in
the integral closure of $S^{-1}A$ in $S^{-1}B$. Conversely, suppose
that $x/f \in S^{-1}B$ is integral over $S^{-1}A$. Then we have
$$
(x/f)^d + \sum\nolimits_{i = 1, \ldots, d} (a_i/f_i) (x/f)^{d - i} = 0
$$
in $S^{-1}B$ for some $a_i \in A$ and $f_i \in S$. This means that
$$
(f'f_1 \ldots f_d x)^d +
\sum\nolimits_{i = 1, \ldots, d}
f^i(f')^if_1^i \ldots f_i^{i - 1} \ldots f_d^i a_i
(f'f_1 \ldots f_dx)^{d - i} = 0
$$
for a suitable $f' \in S$. Hence $f'f_1\ldots f_dx \in B'$ and thus
$x/f \in S^{-1}B'$ as desired.
\end{proof}

\begin{lemma}
\label{lemma-four-rings}
Let
$$
\xymatrix{
S \ar[r] & S' & &
\mathfrak q \ar@{-}[r] & \mathfrak q' \\
R \ar[u] \ar[r] &  R' \ar[u] & &
\mathfrak p \ar@{-}[r] \ar@{-}[u] & \mathfrak p' \ar@{-}[u]
}
$$
be a commutative diagram of rings with primes as indicated.
Assume $R \to S$ is of finite type, and $S \otimes_R R' \to S'$ surjective.
If $R \to S$ is quasi-finite at $\mathfrak q$, then
$R' \to S'$ is quasi-finite at $\mathfrak q'$.
\end{lemma}

\begin{proof}
Write $S \otimes_R \kappa(\mathfrak p) = S_1 \times S_2$
with $S_1$ finite over $\kappa(\mathfrak p)$ and such that
$\mathfrak q$ corresponds to a point of $S_1$ as in
Lemma \ref{lemma-isolated-point}. This product decomposition
induces a corresponding product decomposition for any
$S \otimes_R \kappa(\mathfrak p)$-algebra. In particular,
we obtain $S' \otimes_{R'} \kappa(\mathfrak p') = S'_1 \times S'_2$.
Because $S \otimes_R R' \to S'$ is surjective the canonical map
$(S \otimes_R \kappa(\mathfrak p)) \otimes_{\kappa(\mathfrak p)}
\kappa(\mathfrak p') \to S' \otimes_{R'} \kappa(\mathfrak p')$
is surjective and hence
$S_i \otimes_{\kappa(\mathfrak p)} \kappa(\mathfrak p') \to S'_i$
is surjective. It follows that $S'_1$ is finite over $\kappa(\mathfrak p')$.
The map $S' \otimes_{R'} \kappa(\mathfrak p') \to
\kappa(\mathfrak q')$ factors through $S_1'$
(i.e.\ it annihilates the factor $S_2'$)
because the map $S \otimes_R \kappa(\mathfrak p) \to
\kappa(\mathfrak q)$ factors through $S_1$
(i.e.\ it annihilates the factor $S_2$). Thus
$\mathfrak q'$ corresponds to a point of
$\Spec(S_1')$ in the disjoint union decomposition
of the fibre: $\Spec(S' \otimes_{R'} \kappa(\mathfrak p'))
= \Spec(S_1') \amalg \Spec(S_2')$, see
Lemma \ref{lemma-spec-product}.
Since $S_1'$ is finite over a field, it is an Artinian ring,
and hence $\Spec(S_1')$ is a finite discrete set.
(See Proposition \ref{proposition-dimension-zero-ring}.)
We conclude $\mathfrak q'$ is isolated in its fibre as
desired.
\end{proof}

\begin{lemma}
\label{lemma-make-integral-trick}
Let $R$ be a ring. Let $\varphi : R[x] \to S$ be
a ring map. Let $t \in S$.
Assume that (a) $t$ is integral over $R[x]$,
and (b) there exists a monic $p \in R[x]$ such that
$t \varphi(p) \in \Im(\varphi)$. Then there
exists a $q \in R[x]$ such that $t - \varphi(q)$
is integral over $R$.
\end{lemma}

\begin{proof}
Write $t \varphi(p) = \varphi(r)$ for some $r \in R[x]$.
Using euclidean division, write $r = qp + r'$ with
$q, r' \in R[x]$ and $\deg(r') < \deg(p)$. We may replace
$t$ by $t - \varphi(q)$ which is still integral over
$R[x]$, so that we obtain $t \varphi(p) = \varphi(r')$.
In the ring $S_t$ we may write this as
$\varphi(p) - (1/t) \varphi(r') = 0$.
This implies that $\varphi(x)$ gives an element of the
localization $S_t$ which is integral over
$\varphi(R)[1/t] \subset S_t$. On the other hand,
$t$ is integral over the subring $\varphi(R)[\varphi(x)] \subset S$.
Combined we conclude that $t$ is integral over
the subring $\varphi(R)[1/t] \subset S_t$, see Lemma
\ref{lemma-integral-transitive}. In other words
there exists an equation of the form
$$
t^d + \sum\nolimits_{i < d}
\left(\sum\nolimits_{j = 0, \ldots, n_i} \varphi(r_{i, j})/t^j\right) t^i = 0
$$
in $S_t$ with $r_{i, j} \in R$. This means that
$t^{d + N} +
\sum_{i < d} \sum_{j = 0, \ldots, n_i} \varphi(r_{i, j}) t^{i + N - j} = 0$
in $S$ for some $N$ large enough. In other words
$t$ is integral over $R$.
\end{proof}

\begin{lemma}
\label{lemma-combine-lemmas}
Let $R$ be a ring. Let $\varphi : R[x] \to S$ be
a ring map. Let $t \in S$. Assume $t$ is integral
over $R[x]$. Let $p \in R[x]$, $p = a_0 + a_1x + \ldots +
a_k x^k$ such that $t \varphi(p) \in \Im(\varphi)$.
Then there exists a $q \in R[x]$ and $n \geq 0$
such that $\varphi(a_k)^n t - \varphi(q)$ is integral
over $R$.
\end{lemma}

\begin{proof}
Let $R'$ and $S'$ be the localization of $R$ and $S$ at the element $a_k$.
Let $\varphi' : R'[x] \to S'$ be the localization of $\varphi$.
Let $t' \in S'$ be the image of $t$. Set $p' = p/a_k \in R'[x]$.
Then $t' \varphi'(p') \in \Im(\varphi')$ since $t \varphi(p) \in \Im(\varphi)$.
As $p'$ is monic, by
Lemma \ref{lemma-make-integral-trick} there exists a $q' \in R'[x]$
such that $t' - \varphi'(q')$ is integral over $R'$.
We may choose an $n \geq 0$ and an element $q \in R[x]$
such that $a_k^n q'$ is the image of $q$.
Then $\varphi(a_k)^n t - \varphi(q)$ is an element of $S$ whose
image in $S'$ is integral over $R'$.
By Lemma \ref{lemma-integral-closure-localize}
there exists an $m \geq 0$ such that
$\varphi(a_k)^m(\varphi(a_k)^n t - \varphi(q))$ is integral over $R$.
Thus $\varphi(a_k)^{m + n}t - \varphi(a_k^m q)$
is integral over $R$ as desired.
\end{proof}

\begin{situation}
\label{situation-one-transcendental-element}
Let $R$ be a ring.
Let $\varphi : R[x] \to S$ be finite.
Let
$$
J = \{ g \in S \mid gS \subset \Im(\varphi)\}
$$
be the ``conductor ideal'' of $\varphi$.
Assume that $\varphi(R)$ is integrally closed in $S$.
\end{situation}

\begin{lemma}
\label{lemma-leading-coefficient-in-J}
In Situation \ref{situation-one-transcendental-element}.
Suppose $u \in S$, $a_0, \ldots, a_k \in R$,
$u \varphi(a_0 + a_1x + \ldots + a_k x^k) \in J$.
Then there exists an $m \geq 0$ such that
$u \varphi(a_k)^m \in J$.
\end{lemma}

\begin{proof}
Assume that $S$ is generated by $t_1, \ldots, t_n$
as an $R[x]$-module. In this case
$J = \{ g \in S \mid gt_i \in \Im(\varphi)\text{ for all }i\}$.
Note that each element $u t_i$ is integral over
$R[x]$, see Lemma \ref{lemma-finite-is-integral}.
We have $\varphi(a_0 + a_1x + \ldots + a_k x^k) u t_i \in
\Im(\varphi)$. By Lemma \ref{lemma-combine-lemmas}, for
each $i$ there exists an integer $n_i$ and an element
$q_i \in R[x]$ such that $\varphi(a_k^{n_i}) u t_i - \varphi(q_i)$
is integral over $R$. By assumption this element is in $\varphi(R)$
and hence $\varphi(a_k^{n_i}) u t_i \in \Im(\varphi)$.
It follows that $m = \max\{n_1, \ldots, n_n\}$ works.
\end{proof}

\begin{lemma}
\label{lemma-all-coefficients-in-J}
In Situation \ref{situation-one-transcendental-element}.
Suppose $u \in S$, $a_0, \ldots, a_k \in R$,
$u \varphi(a_0 + a_1x + \ldots + a_k x^k) \in \sqrt{J}$.
Then $u \varphi(a_i) \in \sqrt{J}$ for all $i$.
\end{lemma}

\begin{proof}
Under the assumptions of the lemma we have
$u^n \varphi(a_0 + a_1x + \ldots + a_k x^k)^n \in J$ for
some $n \geq 1$. By Lemma \ref{lemma-leading-coefficient-in-J}
we deduce $u^n \varphi(a_k^{nm}) \in J$ for some $m \geq 1$.
Thus $u \varphi(a_k) \in \sqrt{J}$, and so
$u \varphi(a_0 + a_1x + \ldots + a_k x^k) - u \varphi(a_k x^k) =
u \varphi(a_0 + a_1x + \ldots + a_{k-1} x^{k-1}) \in \sqrt{J}$.
We win by induction on $k$.
\end{proof}

\begin{definition}
\label{definition-strongly-transcendental}
Given an inclusion of rings $R \subset S$ and
an element $x \in S$ we say that $x$ is
{\it strongly transcendental over $R$} if
whenever $u(a_0 + a_1 x + \ldots + a_k x^k) = 0$
with $u \in S$ and $a_i \in R$, then
we have $ua_i = 0$ for all $i$.
\end{definition}

\begin{lemma}
\label{lemma-reduced-strongly-transcendental-minimal-prime}
Suppose $R \subset S$ is an inclusion of reduced rings
and suppose that $x \in S$ is strongly transcendental over $R$.
Let $\mathfrak q \subset S$ be a minimal prime
and let $\mathfrak p = R \cap \mathfrak q$.
Then the image of $x$ in $S/\mathfrak q$ is strongly
transcendental over the subring $R/\mathfrak p$.
\end{lemma}

\begin{proof}
Suppose $u(a_0 + a_1x + \ldots + a_k x^k) \in \mathfrak q$.
By Lemma \ref{lemma-minimal-prime-reduced-ring}
the local ring $S_{\mathfrak q}$ is a field,
and hence $u(a_0 + a_1x + \ldots + a_k x^k) $ is zero
in $S_{\mathfrak q}$. Thus $uu'(a_0 + a_1x + \ldots + a_k x^k) = 0$
for some $u' \in S$, $u' \not\in \mathfrak q$.
Since $x$ is strongly transcendental over $R$ we get
$uu'a_i = 0$ for all $i$. This in turn implies
that $ua_i \in \mathfrak q$.
\end{proof}

\begin{lemma}
\label{lemma-domains-transcendental-not-quasi-finite}
Suppose $R\subset S$ is an inclusion of domains and
let $x \in S$. Assume $x$ is (strongly) transcendental over $R$
and that $S$ is finite over $R[x]$. Then $R \to S$ is not
quasi-finite at any prime of $S$.
\end{lemma}

\begin{proof}
As a first case, assume that $R$ is normal, see
Definition \ref{definition-ring-normal}.
By Lemma \ref{lemma-polynomial-ring-normal}
we see that $R[x]$ is normal.
Take a prime $\mathfrak q \subset S$,
and set $\mathfrak p = R \cap \mathfrak q$.
Assume that the extension $\kappa(\mathfrak p)
\subset \kappa(\mathfrak q)$ is finite.
This would be the case if $R \to S$ is
quasi-finite at $\mathfrak q$.
Let $\mathfrak r = R[x] \cap \mathfrak q$.
Then since $\kappa(\mathfrak p)
\subset \kappa(\mathfrak r) \subset \kappa(\mathfrak q)$
we see that the extension $\kappa(\mathfrak p)
\subset \kappa(\mathfrak r)$ is finite too.
Thus the inclusion $\mathfrak r \supset \mathfrak p R[x]$
is strict. By going down for $R[x] \subset S$,
see Proposition \ref{proposition-going-down-normal-integral},
we find a prime $\mathfrak q' \subset \mathfrak q$,
lying over the prime $\mathfrak pR[x]$. Hence
the fibre $\Spec(S \otimes_R \kappa(\mathfrak p))$
contains a point not equal to $\mathfrak q$,
namely $\mathfrak q'$, whose closure contains $\mathfrak q$ and hence
$\mathfrak q$ is not isolated in its fibre.

\medskip\noindent
If $R$ is not normal, let $R \subset R' \subset K$ be
the integral closure $R'$ of $R$ in its field of fractions
$K$. Let $S \subset S' \subset L$ be the subring $S'$ of
the field of fractions $L$ of $S$ generated by $R'$ and
$S$. Note that by construction the map $S \otimes_R R'
\to S'$ is surjective. This implies that $R'[x] \subset S'$
is finite. Also, the map $S \subset S'$
induces a surjection on $\Spec$, see
Lemma \ref{lemma-integral-overring-surjective}.
We conclude by Lemma \ref{lemma-four-rings} and the normal case
we just discussed.
\end{proof}

\begin{lemma}
\label{lemma-reduced-strongly-transcendental-not-quasi-finite}
Suppose $R \subset S$ is an inclusion of reduced rings.
Assume $x \in S$ is strongly transcendental over $R$,
and $S$ finite over $R[x]$. Then $R \to S$ is not
quasi-finite at any prime of $S$.
\end{lemma}

\begin{proof}
Let $\mathfrak q \subset S$ be any prime.
Choose a minimal prime $\mathfrak q' \subset \mathfrak q$.
According to Lemmas
\ref{lemma-reduced-strongly-transcendental-minimal-prime} and
\ref{lemma-domains-transcendental-not-quasi-finite}
the extension $R/(R \cap \mathfrak q') \subset
S/\mathfrak q'$ is not quasi-finite at the prime corresponding
to $\mathfrak q$. By Lemma \ref{lemma-four-rings}
the extension $R \to S$ is not quasi-finite
at $\mathfrak q$.
\end{proof}
